% Generated draft: requires mathematical and source review.
\documentclass[11pt]{article}
\usepackage[margin=25mm]{geometry}
\usepackage{fontspec}
\setmainfont{Latin Modern Roman}
\usepackage{amsmath,amssymb,mathtools,unicode-math}
\setmathfont{Latin Modern Math}
\usepackage{enumitem,needspace,xurl,hyperref}
\setlength{\parindent}{0pt}
\setlength{\parskip}{5pt}
\setlength{\emergencystretch}{3em}
\newcommand{\N}{\mathbb{N}}
\newcommand{\Z}{\mathbb{Z}}
\newcommand{\Q}{\mathbb{Q}}
\newcommand{\R}{\mathbb{R}}
\newcommand{\C}{\mathbb{C}}
\newcommand{\F}{\mathbb{F}}
\newcommand{\A}{\mathbb{A}}
\newcommand{\PP}{\mathbb{P}}
\newcommand{\E}{\mathbb{E}}
\newcommand{\eps}{\varepsilon}
\newcommand{\dd}{\,\mathrm{d}}
\begin{document}
\section*{UnsolvedMath ID: 30003982}
\textbf{Source problem code:} catalogue entry
\subsection{Definitions and mathematical statement}
For a conventional frame-theoretic interpretation, let \(\mathbb K\in\{\R,\C\}\), let \(n\ge d\ge1\) be integers, and write \(\Phi=(\phi_1,\ldots,\phi_n)\in\mathbb K^{d\times n}\). These choices describe a possible setting; the supplied record does not specify the scalar field or the number of frame vectors. Write \(\Phi^*\) for the adjoint. A Parseval frame satisfies \(\Phi\Phi^*=I_d\), and an equal-norm Parseval frame additionally satisfies \(\|\phi_i\|^2=d/n\) for every \(i\). Define
\[
\mathcal E_{d,n}(\mathbb K)=\left\{\Psi=(\psi_i)_{i=1}^n:\ \Psi\Psi^*=I_d,\quad \|\psi_i\|^2=\frac dn\ \text{for all }i\right\}.
\]
The squared Hilbert--Schmidt distance between ordered frames is
\[
\|\Phi-\Psi\|_{\mathrm{HS}}^2=\sum_{i=1}^n\|\phi_i-\psi_i\|^2.
\]
The frame potential is
\[
\operatorname{FP}(\Phi)=\sum_{i,j=1}^n|\langle\phi_i,\phi_j\rangle|^2
=\operatorname{tr}\bigl((\Phi\Phi^*)^2\bigr).
\]
The supplied record does not explicitly identify its notation \(\|\cdot\|_2\) with the Hilbert--Schmidt norm.

The clean statement asks whether the gradient flow of the frame potential can prove a Paulsen estimate of the form
\[
\|\Phi-\Psi\|_2^2\le C\varepsilon d^{2-\alpha},\qquad 0<\alpha<1.
\]
To display the missing quantifiers without inventing hypotheses, let \(\mathcal A_{d,n}^{\mathbb K}(\varepsilon)\) denote the intended class of approximately equal-norm, approximately Parseval input frames, and let \(I_{d,n}^{\mathbb K}\subset(0,\infty)\) denote the intended range of \(\varepsilon\). Neither this class nor this range is defined in the supplied record. Conditional on specifying them and interpreting the distance as Hilbert--Schmidt distance, a natural uniform formulation is: do there exist constants \(C>0\) and \(\alpha\in(0,1)\), independent of \(d,n,\varepsilon,\Phi\), such that
\[
\begin{gathered}
\forall d,n\in\N\text{ with }n\ge d,\quad
\forall\varepsilon\in I_{d,n}^{\mathbb K},\quad
\forall\Phi\in\mathcal A_{d,n}^{\mathbb K}(\varepsilon),\\
\exists\Psi\in\mathcal E_{d,n}(\mathbb K):\qquad
\|\Phi-\Psi\|_{\mathrm{HS}}^2\le C\varepsilon d^{2-\alpha},
\end{gathered}
\]
with the estimate established using a gradient flow of \(\operatorname{FP}\)? This is a conditional formulation, not a recovery of missing source hypotheses. The field, admissible input class, error range, and permitted dependence of the constants require confirmation. The flow's constraint manifold, metric, normalization, and rule for obtaining \(\Psi\) are also unspecified; an existence estimate alone does not express the requested method of proof.

There is a substantive conflict within the record. Its original statement instead asks for
\[
\|\Phi-\Psi\|_2\le C\varepsilon^{d^2-\alpha},\qquad 0\le\alpha<1.
\]
This changes the power of the distance, the dependence on \(\varepsilon\) and \(d\), and the admissible endpoint for \(\alpha\). The registry labels the clean statement as corrected and verified, but that supplied label does not independently resolve the discrepancy here. Both versions are retained for human review.
\subsection{Short English statement}
The clean record asks whether frame-potential gradient flow can produce an equal-norm Parseval frame with squared distance at most \(C\varepsilon d^{2-\alpha}\), for some \(0<\alpha<1\). The input approximation conditions, norm, scalar field, parameter range, and flow are unspecified. The original statement gives a materially different inequality, so the formulation remains conditional and requires mathematical review.
\subsection{Sources}
\begin{itemize}
\item Supplied local UnsolvedMath registry record, permanent ID \(30003982\), title ``Sharper Paulsen Bounds from Frame Potential Flows,'' source problem code \texttt{catalogue entry}. The record gives proposed year 2018 and names no individual proposer. It labels the clean statement \texttt{corrected\_verified} and reports direct checking of the official report; these are supplied verification claims, not checks performed for this draft.
\item Supplied source citation: \emph{Mini-Workshop: Algebraic, Geometric, and Combinatorial Methods in Frame Theory} (2019), Oberwolfach Reports, pp.~2769--2803, DOI 10.4171/owr/2018/46. \url{https://doi.org/10.4171/owr/2018/46}. The background describes Oberwolfach Reports as records of mathematical workshops and their open-problem discussions. The citation's year 2019 is retained alongside the record's 2018 proposed year and report designation 46/2018; their relationship requires confirmation.
\item Supplied literature reference: \emph{A Proof of the Paulsen Problem}. \url{https://doi.org/10.4230/LIPIcs.ITCS.2019.41}.
\item Supplied literature reference: \emph{Geometric Perspective on Concentration Phenomena in Frame Theory}. \url{https://arxiv.org/abs/2605.03867}.
\item Supplied literature reference: \emph{Oberwolfach Report 46/2018}. \url{https://doi.org/10.4171/owr/2018/46}.
\item Literature assessment dated exactly \texttt{2026-08-21}: the supplied curation labels the problem partially solved and states that the qualitative Paulsen problem is solved, that a deterministic squared-distance bound of order \(\varepsilon d^2\) is known, and that the stronger exponent improvement is unresolved in general. This is supplied dated curation derived from the OWR record, not an independently verified literature assessment. Its time-sensitive status and cited claims require independent verification.
\end{itemize}
\end{document}
